\subsection{机器翻译}
\begin{table}[t]
\begin{center}
\caption{Transformer 在英-德和英-法 newstest2014 测试上的 BLEU 分数超过以前的最先进模型，同时训练成本仅为其一小部分。}
\label{tab:wmt-results}
%\scalebox{1.0}{
\begin{tabular}{lccccc}
\toprule
\multirow{2}{*}{\vspace{-2mm}Model} & \multicolumn{2}{c}{BLEU} & & \multicolumn{2}{c}{Training Cost (FLOPs)} \\
\cmidrule{2-3} \cmidrule{5-6} 
& EN-DE & EN-FR & & EN-DE & EN-FR \\ 
\hline
ByteNet \citep{NalBytenet2017} & 23.75 & & & &\\
Deep-Att + PosUnk \citep{DBLP:journals/corr/ZhouCWLX16} & & 39.2 & & & $1.0\cdot10^{20}$ \\
GNMT + RL \citep{wu2016google} & 24.6 & 39.92 & & $2.3\cdot10^{19}$  & $1.4\cdot10^{20}$\\
ConvS2S \citep{JonasFaceNet2017} & 25.16 & 40.46 & & $9.6\cdot10^{18}$ & $1.5\cdot10^{20}$\\
MoE \citep{shazeer2017outrageously} & 26.03 & 40.56 & & $2.0\cdot10^{19}$ & $1.2\cdot10^{20}$ \\
\hline
\rule{0pt}{2.0ex}Deep-Att + PosUnk Ensemble \citep{DBLP:journals/corr/ZhouCWLX16} & & 40.4 & & &
 $8.0\cdot10^{20}$ \\
GNMT + RL Ensemble \citep{wu2016google} & 26.30 & 41.16 & & $1.8\cdot10^{20}$  & $1.1\cdot10^{21}$\\
ConvS2S Ensemble \citep{JonasFaceNet2017} & 26.36 & \textbf{41.29} & & $7.7\cdot10^{19}$ & $1.2\cdot10^{21}$\\
\specialrule{1pt}{-1pt}{0pt}
\rule{0pt}{2.2ex}Transformer (base model) & 27.3 & 38.1 & & \multicolumn{2}{c}{\boldmath$3.3\cdot10^{18}$}\\
Transformer (big) & \textbf{28.4} & \textbf{41.8} & & \multicolumn{2}{c}{$2.3\cdot10^{19}$} \\
%\hline
%\specialrule{1pt}{-1pt}{0pt}
%\rule{0pt}{2.0ex}
\bottomrule
\end{tabular}
%}
\end{center}
\end{table}


在 WMT 2014 英-德翻译任务中，大型 Transformer 模型（表 \ref{tab:wmt-results} 中的 Transformer (big)）比之前报告的最佳模型（包括集合）提高了超过 $2.0$ 的 BLEU 分数，建立了新的最先进 BLEU 分数 $28.4$。该模型的配置列在表 \ref{tab:variations} 的底行。训练在 8 块 P100 GPU 上耗时 3.5 天。即使是我们的基础模型也超过了所有以前发布的模型和集合，训练成本仅为任何竞争模型的一小部分。

在 WMT 2014 英-法翻译任务中，我们的大型模型达到了 $41.0$ 的 BLEU 分数，超过了所有以前发布的单一模型，训练成本不到以前最先进模型的 1/4。针对英-法翻译训练的 Transformer (big) 模型使用的 dropout 比率为 $P_{drop}=0.1$，而不是 $0.3$。

对于基础模型，我们使用了通过平均最后 5 个检查点获得的单一模型，这些检查点每 10 分钟写入一次。对于大型模型，我们平均了最后 20 个检查点。我们使用了束宽为 4 的束搜索和长度惩罚 $\alpha=0.6$ \citep{wu2016google}。这些超参数是在开发集上实验选择的。我们将推理期间的最大输出长度设置为输入长度 + 50，但尽可能提前终止 \citep{wu2016google}。

表 \ref{tab:wmt-results} 总结了我们的结果，并将我们的翻译质量和训练成本与文献中的其他模型架构进行了比较。我们通过将训练时间、使用的 GPU 数量和对每个 GPU 的持续单精度浮点计算能力的估计相乘来估计用于训练模型的浮点运算次数 \footnote{我们分别为 K80、K40、M40 和 P100 使用了 2.8、3.7、6.0 和 9.5 TFLOPS 的值。}。
%where we compare against the leading machine translation results in the literature. Even our smaller model, with number of parameters comparable to ConvS2S, outperforms all existing single models, and achieves results close to the best ensemble model.

\subsection{模型变体}

\begin{table}[t]
\caption{Transformer 架构的变化。未列出的值与基础模型相同。所有指标都在英-德翻译开发集 newstest2013 上。列出的困惑度是按照我们的字节对编码的词片段级别，不应与按词级别的困惑度进行比较。}
\label{tab:variations}
\begin{center}
\vspace{-2mm}
%\scalebox{1.0}{
\begin{tabular}{c|ccccccccc|ccc}
\hline\rule{0pt}{2.0ex}
 & \multirow{2}{*}{$N$} & \multirow{2}{*}{$\dmodel$} &
\multirow{2}{*}{$\dff$} & \multirow{2}{*}{$h$} & 
\multirow{2}{*}{$d_k$} & \multirow{2}{*}{$d_v$} & 
\multirow{2}{*}{$P_{drop}$} & \multirow{2}{*}{$\epsilon_{ls}$} &
train & PPL & BLEU & params \\
 & & & & & & & & & steps & (dev) & (dev) & $\times10^6$ \\
% & & & & & & & & & & & & \\
\hline\rule{0pt}{2.0ex}
base & 6 & 512 & 2048 & 8 & 64 & 64 & 0.1 & 0.1 & 100K & 4.92 & 25.8 & 65 \\
\hline\rule{0pt}{2.0ex}
\multirow{4}{*}{(A)}
& & & & 1 & 512 & 512 & & & & 5.29 & 24.9 &  \\
& & & & 4 & 128 & 128 & & & & 5.00 & 25.5 &  \\
& & & & 16 & 32 & 32 & & & & 4.91 & 25.8 &  \\
& & & & 32 & 16 & 16 & & & & 5.01 & 25.4 &  \\
\hline\rule{0pt}{2.0ex}
\multirow{2}{*}{(B)}
& & & & & 16 & & & & & 5.16 & 25.1 & 58 \\
& & & & & 32 & & & & & 5.01 & 25.4 & 60 \\
\hline\rule{0pt}{2.0ex}
\multirow{7}{*}{(C)}
& 2 & & & & & & & &            & 6.11 & 23.7 & 36 \\
& 4 & & & & & & & &            & 5.19 & 25.3 & 50 \\
& 8 & & & & & & & &            & 4.88 & 25.5 & 80 \\
& & 256 & & & 32 & 32 & & &    & 5.75 & 24.5 & 28 \\
& & 1024 & & & 128 & 128 & & & & 4.66 & 26.0 & 168 \\
& & & 1024 & & & & & &         & 5.12 & 25.4 & 53 \\
& & & 4096 & & & & & &         & 4.75 & 26.2 & 90 \\
\hline\rule{0pt}{2.0ex}
\multirow{4}{*}{(D)}
& & & & & & & 0.0 & & & 5.77 & 24.6 &  \\
& & & & & & & 0.2 & & & 4.95 & 25.5 &  \\
& & & & & & & & 0.0 & & 4.67 & 25.3 &  \\
& & & & & & & & 0.2 & & 5.47 & 25.7 &  \\
\hline\rule{0pt}{2.0ex}
(E) & & \multicolumn{7}{c}{positional embedding instead of sinusoids} & & 4.92 & 25.7 & \\
\hline\rule{0pt}{2.0ex}
big & 6 & 1024 & 4096 & 16 & & & 0.3 & & 300K & \textbf{4.33} & \textbf{26.4} & 213 \\
\hline
\end{tabular}
%}
\end{center}
\end{table}


%Table \ref{tab:ende-results}. Our base model for this task uses 6 attention layers, 512 hidden dim, 2048 filter dim, 8 attention heads with both attention and symbol dropout of 0.2 and 0.1 respectively. Increasing the filter size of our feed forward component to 8192 increases the BLEU score on En $\to$ De by $?$. For both the models, we use beam search decoding of size $?$ and length penalty with an alpha of $?$ \cite? \todo{Update results}

为了评估 Transformer 不同组件的重要性，我们以不同的方式改变基础模型，测量英-德翻译在开发集 newstest2013 上的性能变化。我们使用了前一部分所述的束搜索，但不使用检查点平均。我们在表 \ref{tab:variations} 中呈现了这些结果。  

在表 \ref{tab:variations} 的 (A) 行中，我们改变注意力头的数量和注意力 key 与 value 的维度，保持计算量不变，如第 \ref{sec:multihead} 节所述。虽然单头注意力比最佳设置差 0.9 个 BLEU 分数，但质量也会随着头数过多而下降。

在表 \ref{tab:variations} 的 (B) 行中，我们观察到减少注意力 key 大小 $d_k$ 会降低模型质量。这表明确定兼容性并不容易，更复杂的兼容性函数可能比点积更有益。我们进一步在 (C) 和 (D) 行观察到，如预期的那样，更大的模型更好，dropout 在避免过拟合方面非常有帮助。在 (E) 行中，我们用学习的位置嵌入替换了我们的正弦位置编码 \citep{JonasFaceNet2017}，观察到与基础模型几乎相同的结果。

%To evaluate the importance of different components of the Transformer, we use our base model to ablate on a single hyperparameter at each time and measure the change in performance on English$\to$German translation. Our variations in Table~\ref{tab:variations} show that the number of attention layers and attention heads is the most important architecture hyperparamter However, the we do not see performance gains beyond 6 layers, suggesting that we either don't have enough data to train a large model or we need to turn up regularization. We leave this exploration for future work. Among our regularizers, attention dropout has the most significant impact on performance.


%Increasing the width of our feed forward component helps both on log ppl and Accuracy \marginpar{Intuition?}
%Using dropout to regularize our models helps to prevent overfitting

\subsection{英文句法成分解析}

\begin{table}[t]
\begin{center}
\caption{Transformer 在英文句法成分解析中泛化良好（结果基于 WSJ 的第 23 部分）}
\label{tab:parsing-results}
%\scalebox{1.0}{
\begin{tabular}{c|c|c}
\hline
{\bf Parser}  & {\bf Training} & {\bf WSJ 23 F1} \\ \hline
Vinyals \& Kaiser el al. (2014) \cite{KVparse15}
  & WSJ only, discriminative & 88.3 \\
Petrov et al. (2006) \cite{petrov-EtAl:2006:ACL}
  & WSJ only, discriminative & 90.4 \\
Zhu et al. (2013) \cite{zhu-EtAl:2013:ACL}
  & WSJ only, discriminative & 90.4   \\
Dyer et al. (2016) \cite{dyer-rnng:16}
  & WSJ only, discriminative & 91.7   \\
\specialrule{1pt}{-1pt}{0pt}
Transformer (4 layers)  &  WSJ only, discriminative & 91.3 \\
\specialrule{1pt}{-1pt}{0pt}   
Zhu et al. (2013) \cite{zhu-EtAl:2013:ACL}
  & semi-supervised & 91.3 \\
Huang \& Harper (2009) \cite{huang-harper:2009:EMNLP}
  & semi-supervised & 91.3 \\
McClosky et al. (2006) \cite{mcclosky-etAl:2006:NAACL}
  & semi-supervised & 92.1 \\
Vinyals \& Kaiser el al. (2014) \cite{KVparse15}
  & semi-supervised & 92.1 \\
\specialrule{1pt}{-1pt}{0pt}
Transformer (4 layers)  & semi-supervised & 92.7 \\
\specialrule{1pt}{-1pt}{0pt}   
Luong et al. (2015) \cite{multiseq2seq}
  & multi-task & 93.0   \\
Dyer et al. (2016) \cite{dyer-rnng:16}
  & generative & 93.3   \\
\hline
\end{tabular}
\end{center}
\end{table}

为了评估 Transformer 是否可以泛化到其他任务，我们在英文句法成分解析上进行了实验。这个任务呈现了特定的挑战：输出受到强结构约束，长度显著大于输入。此外，RNN 序列到序列模型在小数据体系中无法获得最先进的结果 \cite{KVparse15}。

我们在 Penn Treebank 的华尔街日报 (WSJ) 部分上训练了一个 4 层的 Transformer，其中 $d_{model} = 1024$，约 40K 个训练句子 \citep{marcus1993building}。我们也在半监督设置中训练它，使用了更大的高信度和 BerkleyParser 语料库，约包含 1700 万个句子 \citep{KVparse15}。对于仅 WSJ 设置，我们使用了 16K 个 token 的词汇表，对于半监督设置，使用了 32K 个 token 的词汇表。

我们只进行了少量实验来选择 dropout（包括注意力和残差 dropout，见第 \ref{sec:reg} 部分）、学习率和第 22 部分开发集上的束大小，所有其他参数保持与英-德基础翻译模型相同。在推理期间，我们将最大输出长度增加到输入长度 + 300。对于仅 WSJ 和半监督设置，我们都使用了束大小 21 和 $\alpha=0.3$。

表 \ref{tab:parsing-results} 中的结果显示，尽管缺乏任务特定的调优，我们的模型表现得出奇地好，除了循环神经网络语法 \cite{dyer-rnng:16} 之外，优于所有以前报告的模型。

与 RNN 序列到序列模型相反 \citep{KVparse15}，Transformer 甚至在仅用 40K 个句子的 WSJ 训练集进行训练时，也优于 BerkeleyParser \cite{petrov-EtAl:2006:ACL}。
